\section*{Hoofdstuk \ref{ch:g12:cells-and-electrolysis} --- Elektrochemische cellen en elektrolyse}

\begin{solution}{exo:g12:cells-and-electrolysis:1}
De zinkelektrode is de \omterm{def:g12:cells-and-electrolysis:anode}{anode}, \ce{Zn -> Zn^{2+} + 2e-}; de koperelektrode
is de \omterm{def:g12:cells-and-electrolysis:anode}{kathode}, \ce{Cu^{2+} + 2e- -> Cu}.
\end{solution}

\begin{solution}{exo:g12:cells-and-electrolysis:2}
% ledger: const:F
$Q = 0.20 \times 1800 = \qty{360}{C}$; $n(e^-) = 360/96485 =
\qty{3.7e-3}{mol}$.
\end{solution}

\begin{solution}{exo:g12:cells-and-electrolysis:3}
Ze sluit de stroomkring en houdt elke oplossing elektrisch neutraal,
doordat haar \omterm{def:g9:ions:ion}{ionen} de compartimenten in bewegen. Zonder \omterm{def:g12:cells-and-electrolysis:cell}{zoutbrug} is de
stroomkring open: er loopt geen stroom, de voltmeter wijst nul aan en een
lampje blijft uit.
\end{solution}

\begin{solution}{exo:g12:cells-and-electrolysis:4}
$2.5 \times 3600 = \qty{9000}{C}$.
\end{solution}

\begin{solution}{exo:g12:cells-and-electrolysis:5}
Nee: de reactie verloopt in de richting tegengesteld aan de spontane. De
spanningsbron levert de energie.
\end{solution}

\begin{solution}{exo:g12:cells-and-electrolysis:6}
\omterm{def:g12:cells-and-electrolysis:anode}{Anode} (koper, $-$): \ce{Cu -> Cu^{2+} + 2e-}. \omterm{def:g12:cells-and-electrolysis:anode}{Kathode} (zilver, $+$):
\ce{Ag+ + e- -> Ag}. Totaal: \ce{Cu + 2Ag+ -> Cu^{2+} + 2Ag}.
\end{solution}

\begin{solution}{exo:g12:cells-and-electrolysis:7}
Drie jaar is $3 \times 365 \times 24 = \qty{26280}{h}$;
$0.010 \times 26280 = \qty{263}{mA.h}$.
\end{solution}

\begin{solution}{exo:g12:cells-and-electrolysis:8}
% ledger: const:F, aw:Ag
$Q = 0.50 \times 2400 = \qty{1200}{C}$; $n(e^-) = n(\ce{Ag}) =
1200/96485 = \qty{0.0124}{mol}$; $m = 0.0124 \times 107.9 =
\qty{1.34}{g}$.
\end{solution}

\begin{solution}{exo:g12:cells-and-electrolysis:9}
De stroom loopt in de draad van het koper ($+$) naar het zink ($-$),
tegengesteld aan de \omterm{def:g9:inside-the-atom:nucleus}{elektronen}. De \ce{K+}-ionen bewegen naar het
kopercompartiment, dat \ce{Cu^{2+}} verliest; de \ce{NO3-}-ionen naar het
zinkcompartiment, dat \ce{Zn^{2+}} wint.
\end{solution}

\begin{solution}{exo:g12:cells-and-electrolysis:10}
In beide gevallen is de \omterm{def:g12:cells-and-electrolysis:anode}{anode} de plek waar de \omterm{def:g11:redox:oxidation}{oxidatie} plaatsvindt. In
een cel is de \omterm{def:g12:cells-and-electrolysis:anode}{anode} $-$, verloopt de reactie vanzelf en wordt chemische
energie elektrische energie; bij een \omterm{def:g12:cells-and-electrolysis:electrolysis}{elektrolyse} is de \omterm{def:g12:cells-and-electrolysis:anode}{anode} verbonden
met de $+$ van de spanningsbron, wordt de reactie afgedwongen en wordt
elektrische energie chemische energie.
\end{solution}

\begin{solution}{exo:g12:cells-and-electrolysis:11}
\qty{12}{mL}: de vergelijking \ce{2H2O -> 2H2 + O2} geeft twee \omterm{def:g7:atoms-and-molecules:molecule}{moleculen}
waterstof op één \omterm{def:g7:atoms-and-molecules:molecule}{molecuul} zuurstof, en gelijke hoeveelheden gas nemen
gelijke volumes in.
\end{solution}

\begin{solution}{exo:g12:cells-and-electrolysis:12}
% ledger: aw:Zn, const:F
$n(\ce{Zn}) = 5.0/65.4 = \qty{0.076}{mol}$;
$n(\ce{Cu^{2+}}) = 0.100 \times 0.50 = \qty{0.050}{mol}$; ze reageren één
op één, dus de koperionen raken het eerst op. $n(e^-) = 2 \times 0.050
= \qty{0.100}{mol}$, $Q = \qty{9.65e3}{C} = \qty{2.68}{A.h}$. Bij
\qty{50}{mA}: $2.68/0.050 = \qty{54}{h}$.
\end{solution}

\begin{solution}{exo:g12:cells-and-electrolysis:13}
$Q = 2.0 \times 3600 = \qty{7200}{C}$; $E = 7200 \times 1.5 =
\qty{1.08e4}{J}$, dat is $1.5 \times 2.0 = \qty{3.0}{W.h}$.
\end{solution}

\begin{solution}{exo:g12:cells-and-electrolysis:14}
% ledger: const:F
$Q = \qty{7200}{C}$, $n(e^-) = \qty{0.0746}{mol}$. Twee \omterm{def:g9:inside-the-atom:nucleus}{elektronen} per
\ce{H2}: $n(\ce{H2}) = \qty{0.0373}{mol}$, $V = \qty{0.90}{L}$. Vier per
\ce{O2}: $n(\ce{O2}) = \qty{0.0187}{mol}$, $V = \qty{0.45}{L}$.
\end{solution}

\begin{solution}{exo:g12:cells-and-electrolysis:15}
Het koper dat aan de \omterm{def:g12:cells-and-electrolysis:anode}{anode} \omterm{def:g2:mixing-and-dissolving:dissolve}{oplost}, is gelijk aan het koper dat neerslaat,
omdat aan beide kanten dezelfde \omterm{def:g9:inside-the-atom:nucleus}{elektronen} passeren:
$2.84/3.00 = \qty{94.7}{\percent}$ koper.
\end{solution}

\begin{solution}{pb:g12:cells-and-electrolysis:1}
% ledger: const:F, const:NA, aw:Li, dch:CH4
\textbf{1.} De \omterm{def:g12:cells-and-electrolysis:anode}{anode}: daar wordt lithium geoxideerd.

\textbf{2.} De grafietelektrode, de \omterm{def:g12:cells-and-electrolysis:anode}{anode}.

\textbf{3.} De \omterm{def:g9:inside-the-atom:nucleus}{elektronen} gaan door de telefoon van de grafietelektrode
naar de oxide-elektrode; de lithiumionen steken de elektrolyt in dezelfde
richting over, van grafiet naar oxide.

\textbf{4.} Haar reactie kan door een lader achteruit gedreven worden.

\textbf{5.} \qty{4.000}{A.h}, dat is $4.000 \times 3600 =
\qty{14400}{C}$.

\textbf{6.} $14400/96485 = \qty{0.1492}{mol}$.

\textbf{7.} $0.1492 \times \num{6.022e23} = \num{8.99e22}$ \omterm{def:g9:inside-the-atom:nucleus}{elektronen}.

\textbf{8.} $4.000/0.25 = \qty{16}{h}$.

\textbf{9.} $0.20 \times 14400 = \qty{2880}{C}$.

\textbf{10.} $E = 14400 \times 3.7 = \qty{5.3e4}{J}$, ongeveer
\qty{53}{kJ}.

\textbf{11.} $53280/3600 = \qty{14.8}{W.h}$.

\textbf{12.} De \omterm{def:g4:what-a-fire-needs:burning}{verbranding} van \qty{1}{g} methaan levert iets meer
energie (\qty{55.7}{kJ}) dan een volle telefoonaccu opslaat.

\textbf{13.} $1000/14.8 = 67.6$: 67 volledige oplaadbeurten.

\textbf{14.} \ce{Li+ + e- -> Li}: een \omterm{def:g11:redox:oxidation}{reductie}.

\textbf{15.} De lader dwingt de reactie in de richting tegengesteld aan
de spontane.

\textbf{16.} $14400/2.0 = \qty{7200}{s}$, \qty{2.0}{h}.

\textbf{17.} Eén lithiumion per \omterm{def:g9:inside-the-atom:nucleus}{elektron}: \qty{0.1492}{mol}.

\textbf{18.} $0.1492 \times 6.9 = \qty{1.03}{g}$.

\textbf{19.} Bij één volledige oplaadbeurt pendelt er ongeveer
\qty{1.03}{g} lithium tussen de elektroden.
\end{solution}
